\clearpage
\section{Arithmetic coherence and the remaining machine cases}
\label{sec:copying-arithmetic-coherence}
The original arithmetic rules now have complete successor-state coherence
under every constructor of the actual contextual rewrite relation. This covers
unary natural negation, addition and truncated subtraction. A second result
exhausts the computation alphabet: any pair of successors is recursively
isomorphic, or the original state requests application or copying. At that milestone these were
the two remaining computation families for successor coherence.
Section~\ref{sec:copying-application-coherence} subsequently completes application,
leaving only copying in this classification. This is not yet a proof of
full-machine determinism, refocusing or robustness.

Arithmetic differs from the completed store comparisons because it removes
both operand edges and their internal wires. The proof records the surviving
occurrences in the original source, derives this record from every actual
rewrite, and compares two results through their common source. No deletion
certificate is added to the operational premises. The binary rule remains
Figure~10(a), and unary negation remains the agreed natural-number operation
$-n=0-n=0$. The terms of Proposition~12.2, Figure~16's translation, and the
stability conditions of other operations are unchanged.

Every successful unary redex execution has exactly five transitions, and every
successful binary redex execution has exactly seven, even after arbitrary finite
depth padding. These statements quantify over all successful runs of the full
machine. A binary redex attached below an arbitrary host returns control in
seven transitions, after which the host's own successful cost is retained.
The nested arithmetic regression has exactly eleven successful transitions.
These are explicit graph redexes, not simplified replacements for complete
surface-term translations.

The dependency chain is: unique source recognition; survivor certificates;
renaming and context transport; certificates for the original local rules;
induction over actual rewrites; whole-state coherence; exact cost laws; and
exhaustive classification of the remaining heads. Appendix~B.1 removes
arithmetic from its outstanding successor comparisons. Its four broad work
packages remain open, and the mathematical-effort estimate stays at roughly
50--65\%, with low confidence.

\Needspace{12\baselineskip}
\subsection{Surviving vertices and edges}

\termdef{internal-arithmetic-targets}
A head has \emph{internal arithmetic targets} when each port in its target
list is an internal vertex of the source graph. This property prevents
arithmetic deletion from removing a host boundary wire after insertion. It is
proved for the original redexes and transported through all rewrite contexts;
it does not assume global acyclicity or injectivity of boundary attachment.

\lean{Copying/ArithmeticReduction}{9}{10}{Arithmetic.InternalTargets}

\termdef{arithmetic-deletion-certificate}
Fix a source graph $G$, its arithmetic head $e$, and the numeric result $r$.
A \emph{source-relative deletion certificate} for a successor $H$ consists of
injective maps $j_V:V_H\to V_G$ on internal vertices and
$j_E:E_H\to E_G$ on edges, with exact images
\[
\begin{split}
\operatorname{im}(j_V)&=\{x\mid\operatorname{internal}(x)\notin\operatorname{tgt}_G(e)\},\\
\operatorname{im}(j_E)&=\{a\mid\neg\operatorname{Adjacent}_G(e,a)\}.
\end{split}
\]
Here adjacency means that a target of $e$ is a source of $a$. The port map
$\widehat j_V$ fixes external inputs and outputs and acts by $j_V$ internally.
Internal labels and every ordered source list are preserved. An edge mapped
to $e$ has numeric label $r$ and empty target list. Every other edge has a
recursively equivalent label and the same ordered targets after mapping by
$\widehat j_V$. Thus the certificate describes the entire successor, including
all retained boxes; it is stronger than identifying the numeric result alone.
The local redex proofs below show that the head survives and precisely its
numeric operands and operand wires disappear.

\lean{Copying/ArithmeticReduction}{14}{26}{Arithmetic.Reduction}

Two injective maps with the same image determine a bijection between their
domains: send an element to the unique element with the same image. Classical
choice selects the preimage; injectivity proves both inverse equations. The
commuting equation says that applying the second embedding after this bijection
recovers the first embedding. Applying this construction separately to $j_V$
and $j_E$ gives the vertex and edge bijections between two certified successors.
Their commuting equations follow directly; the port equation adds the input
and output cases, where both maps are the identity.

\lean{Copying/ArithmeticReduction}{32}{40}{Arithmetic.Reduction.sameRange}

\lean{Copying/ArithmeticReduction}{42}{45}{Arithmetic.Reduction.sameRange_commute}

\lean{Copying/ArithmeticReduction}{47}{50}{Arithmetic.Reduction.vertexBijection}

\lean{Copying/ArithmeticReduction}{52}{55}{Arithmetic.Reduction.edgeBijection}

\lean{Copying/ArithmeticReduction}{57}{58}{Arithmetic.Reduction.vertex_commute}

\lean{Copying/ArithmeticReduction}{60}{61}{Arithmetic.Reduction.edge_commute}

\lean{Copying/ArithmeticReduction}{63}{65}{Arithmetic.Reduction.port_commute}

These bijections form a recursive graph isomorphism. Internal labels agree
through their common source labels. At the arithmetic head both edge labels
are the same numeric constant; elsewhere compose the two recursive label
comparisons through $G$. For incidences, map both candidate lists into $G$,
use their certificate equations and the commuting equations, then cancel the
injective port map. The head's two target lists are both empty.
If the two resumed tokens also map to the same source port, the port commuting
equation and injectivity identify them. Both states are in search mode, so this
upgrades the graph isomorphism to a complete state isomorphism.

\lean{Copying/ArithmeticReduction}{68}{93}{Arithmetic.Reduction.iso}

\lean{Copying/ArithmeticReduction}{96}{101}{Arithmetic.Reduction.stateIso}

\Needspace{12\baselineskip}
\subsection{Renaming the source and successor}

Port mapping respects composition of internal-vertex maps, by the three
port constructors. Mapping an entire list respects the same composition, by
list-map composition and the pointwise port equation. These elementary laws
control all incidence equations under renaming.

\lean{Copying/ArithmeticReductionRenaming}{9}{10}{Arithmetic.Reduction.port_comp}

\lean{Copying/ArithmeticReductionRenaming}{12}{15}{Arithmetic.Reduction.map_ports_comp}

To re-present the successor by an isomorphism $K\cong H$, compose its carrier
bijections with $H$'s maps into $G$. Surjectivity leaves the images unchanged,
and composition preserves injectivity. Recursive label comparison transports
the numeric head label exactly and all other labels relationally. The incidence
laws compose; an empty mapped target list implies the original list is empty.
Thus the certificate is preserved under arbitrary recursive successor renaming.

\lean{Copying/ArithmeticReductionRenaming}{18}{47}{Arithmetic.Reduction.targetRename}

To re-present the source by $f:G\cong N$, compose the certificate maps with
$f$ and replace the selected head by $f(e)$. Target membership reflects through
the inverse vertex bijection, and adjacency reflects through the inverse edge
bijection. These facts transfer the exact-image conditions. Labels and incidences
then compose with the corresponding isomorphism laws. The two port-transport
lemmas state the resulting composition equations explicitly, so the resumed
token can subsequently be transported without rebuilding this argument.

\lean{Copying/ArithmeticReductionRenaming}{50}{92}{Arithmetic.Reduction.sourceRename}

\lean{Copying/ArithmeticReductionRenaming}{94}{95}{Arithmetic.Reduction.targetRename_port}

\lean{Copying/ArithmeticReductionRenaming}{97}{98}{Arithmetic.Reduction.sourceRename_port}

\Needspace{12\baselineskip}
\subsection{Deleting an arithmetic redex inside a plugged graph}

Let $C[h:=G]$ be a typed plugging. An inserted component port can equal an
internal vertex of that component only if it was already that internal vertex:
inputs and outputs attach to host ports, whose carrier tags are disjoint from
component-internal tags. When the selected head has internal targets, its
plugged target vertices are exactly the inserted original target vertices.
It therefore cannot be adjacent to a retained host edge, and adjacency to an
inserted edge is equivalent to adjacency inside $G$. The forward direction
reflects the shared internal port using the preceding disjointness fact; the
reverse direction uses incidence transport. Internal-target closure itself
survives plugging because insertion preserves component-internal vertices.

\lean{Copying/ArithmeticReductionPlug}{13}{18}{Arithmetic.Reduction.inserted_internal}

\lean{Copying/ArithmeticReductionPlug}{20}{29}{Arithmetic.Reduction.plug_target}

\lean{Copying/ArithmeticReductionPlug}{31}{37}{Arithmetic.Reduction.plug_adjacent_host}

\lean{Copying/ArithmeticReductionPlug}{39}{47}{Arithmetic.Reduction.plug_adjacent_inner}

\lean{Copying/ArithmeticReductionPlug}{49}{53}{Arithmetic.Reduction.plug_closed}

The whole successor's vertex map is the identity on host vertices and the
component certificate map on inserted vertices. Mapping an inserted port this
way agrees with first mapping it in the component and then inserting it into
the original plugged graph. Internal ports are immediate; boundary ports
follow by considering the host port to which they attach. The host-port equation
is simply the identity in each of the three port cases.

\lean{Copying/ArithmeticReductionPlug}{55}{65}{Arithmetic.Reduction.plug_port}

\lean{Copying/ArithmeticReductionPlug}{67}{69}{Arithmetic.Reduction.plug_host}

Define the whole edge map likewise: retain each host edge and use the
component edge map inside the hole. Disjoint carrier tags and component
injectivity prove injectivity. The preceding target and adjacency lemmas give
the exact vertex and edge images: every host occurrence survives, and the
component removes precisely its original operand occurrences. Host labels are
unchanged; component labels use the original certificate. Ordered incidence
lists commute by the two port equations. This constructs a deletion certificate
for the whole plugging, without imposing a common host decomposition on two
independently chosen arithmetic matches.

\lean{Copying/ArithmeticReductionPlug}{72}{164}{Arithmetic.Reduction.plug}

\Needspace{12\baselineskip}
\subsection{Simultaneous expansion and dependent component carriers}

Suppose two filling families $F$ and $K$ agree at every host edge except
$h$, where a certificate compares $F(h)$ with $K(h)$. Equality of two hypernets
induces a recursive isomorphism in the reverse direction by eliminating the
equality. For each component, define its vertex and edge maps to be the
certificate maps when its index is $h$, and the equality-induced isomorphism
maps otherwise. At $h$ these definitions reduce to the certificate maps.
The whole expansion maps preserve the host vertices and every component index,
using these maps within the corresponding dependent component carriers.

\lean{Copying/ArithmeticReductionExpansion}{13}{14}{Arithmetic.Reduction.equalIso}

\lean{Copying/ArithmeticReductionExpansion}{16}{19}{Arithmetic.Reduction.fiberVertices}

\lean{Copying/ArithmeticReductionExpansion}{21}{24}{Arithmetic.Reduction.fiberEdges}

\lean{Copying/ArithmeticReductionExpansion}{26}{27}{Arithmetic.Reduction.fiber_vertices_same}

\lean{Copying/ArithmeticReductionExpansion}{29}{30}{Arithmetic.Reduction.fiber_edges_same}

\lean{Copying/ArithmeticReductionExpansion}{32}{34}{Arithmetic.Reduction.expansionVertices}

\lean{Copying/ArithmeticReductionExpansion}{36}{38}{Arithmetic.Reduction.expansionEdges}

Mapping a port inserted from any component commutes with inserting its
component image. Internal ports retain their component index; boundary ports
attach to the same host port on both sides. A port inserted at index $e$ can
be an internal vertex tagged by $e$ only when it was that internal vertex.
The selected head's expanded targets are therefore exactly its own tagged
internal targets. Adjacency within the selected component reflects and
preserves adjacency there. Adjacency to another component is impossible:
internal ports have different tags, and boundary attachments lie in the host.
In particular the selected head still has only internal targets.

\lean{Copying/ArithmeticReductionExpansion}{40}{50}{Arithmetic.Reduction.expansion_port}

\lean{Copying/ArithmeticReductionExpansion}{52}{61}{Arithmetic.Reduction.expansion_internal}

\lean{Copying/ArithmeticReductionExpansion}{63}{71}{Arithmetic.Reduction.expansion_target}

\lean{Copying/ArithmeticReductionExpansion}{73}{81}{Arithmetic.Reduction.expansion_adjacent_same}

\lean{Copying/ArithmeticReductionExpansion}{83}{97}{Arithmetic.Reduction.expansion_adjacent_other}

\lean{Copying/ArithmeticReductionExpansion}{99}{104}{Arithmetic.Reduction.expansion_closed}

These facts construct the expansion certificate. Each component map is
injective, by the given certificate at $h$ and by an isomorphism elsewhere.
Preserved indices give injectivity of the whole maps. In the selected component
the image conditions are exactly the certificate conditions; every other
component is covered by the inverse isomorphism, and every host vertex survives.
No other component is adjacent to the selected head. Internal labels, recursively
equivalent edge labels, and incidences follow component by component; the port
commutation law handles attachment. This proof also accounts for unchanged
components containing arbitrary boxes or cycles.

\lean{Copying/ArithmeticReductionExpansion}{106}{266}{Arithmetic.Reduction.expand}

\Needspace{12\baselineskip}
\subsection{Certificates for the original numeric rules}

A two-element index is either zero or one, by its finite bound. Enumerating
the serial redex carriers shows that every unary internal vertex is its sole
operand wire and every binary internal vertex is one of its two ordered operand
wires. Primitive numeric edges have no internal vertices of their own. Thus
all internal vertices occur in the head's target list, and every head target
is internal. Enumerating the two unary edges, or the three binary edges, shows
that the head is adjacent exactly to the other edges. The head is not adjacent
to itself, because its source is the external input.

\lean{Copying/ArithmeticLocalReduction}{8}{8}{Arithmetic.fin_two}

\lean{Copying/ArithmeticLocalReduction}{10}{15}{Arithmetic.unary_internal_deleted}

\lean{Copying/ArithmeticLocalReduction}{17}{26}{Arithmetic.binary_internal_deleted}

\lean{Copying/ArithmeticLocalReduction}{28}{30}{Arithmetic.unary_internal_targets}

\lean{Copying/ArithmeticLocalReduction}{32}{36}{Arithmetic.binary_internal_targets}

\lean{Copying/ArithmeticLocalReduction}{38}{48}{Arithmetic.unary_adjacent_iff}

\lean{Copying/ArithmeticLocalReduction}{50}{68}{Arithmetic.binary_adjacent_iff}

A general local construction now suffices. Suppose a closed redex has all
internal vertices targeted by its head, the head is adjacent precisely to all
other edges, and its source list is the singleton external input. The result
corolla has no internal vertices, so its vertex map is the empty map. Map its
unique edge to the old head. The vertex image is empty by the first hypothesis;
the edge image is exactly the head by the second. The singleton source and empty
target list satisfy the incidence requirements, and the result label is numeric
by construction. Applying this construction to unary and binary redexes gives
the original arithmetic rules' certificates, with results $0-n$ and
$\operatorname{eval}(o,m,n)$ respectively.

\lean{Copying/ArithmeticLocalReduction}{71}{98}{Arithmetic.result_reduction}

\lean{Copying/ArithmeticLocalReduction}{100}{102}{Arithmetic.unary_reduction}

\lean{Copying/ArithmeticLocalReduction}{104}{106}{Arithmetic.binary_reduction}

\Needspace{12\baselineskip}
\subsection{Recognized arithmetic requests and actual contextual rewrites}

\termdef{arithmetic-request}
An \emph{arithmetic request} is either unary negation with operand $n$, or a
binary code $o$ with ordered operands $m,n$. Its recognized machine choice has
the same data; its result evaluates the indicated natural-number operation.
The map to machine choices is injective, because the constructors and their
fields distinguish the requests.

\lean{Copying/ArithmeticCoherence}{8}{10}{Arithmetic.Request}

\lean{Copying/ArithmeticCoherence}{12}{14}{Arithmetic.Request.choice}

\lean{Copying/ArithmeticCoherence}{16}{18}{Arithmetic.Request.result}

\lean{Copying/ArithmeticCoherence}{20}{21}{Arithmetic.Request.choice_injective}

A complete arithmetic certificate additionally specifies the source head,
recognition of its primitive label at the source token, internal-target closure,
a deletion certificate, and the equation mapping the resumed successor token
back to the original source token. Internal-target closure is invariant under
recursive isomorphism: reflect target membership, use internality in the source,
then map the internal vertex forward.

\lean{Copying/ArithmeticCoherence}{24}{29}{Arithmetic.Certificate}

\lean{Copying/ArithmeticCoherence}{31}{37}{Arithmetic.internalTargets_rename}

For renaming, transport the head by the match embedding, rename the source
and then the successor of the deletion certificate, and cancel the successor
isomorphism with its inverse in the token equation. For plugging, use the
plugged deletion certificate and head, together with the inserted-port
commutation law. For expansion, use the expanded certificate and the
component-port equation at $h$. In each case the resumed token is tracked through
the same carrier map as the graph. The original unary and binary rules supply
complete certificates using the explicit local reductions; both tokens are the
external input, so their token equation is reflexive.

\lean{Copying/ArithmeticCoherence}{39}{51}{Arithmetic.Certificate.rename}

\lean{Copying/ArithmeticCoherence}{53}{62}{Arithmetic.Certificate.plug}

\lean{Copying/ArithmeticCoherence}{64}{77}{Arithmetic.Certificate.expand}

\lean{Copying/ArithmeticCoherence}{80}{86}{Arithmetic.unary_certificate}

\lean{Copying/ArithmeticCoherence}{88}{94}{Arithmetic.binary_certificate}

Recognition can be reflected into the component of an actual rewrite even
though a general match embedding need not be surjective. The component's actual
rewrite has some recognized choice by the earlier selection theorem. Transport
that choice to the host and compare it with the specified host recognition.
Uniqueness of source recognition identifies the choices, leaving the required
component recognition. The actual-rewrite premise is essential to this argument.

\lean{Copying/ArithmeticCoherence}{101}{105}{Combined.rewrite_recognition_reflect}

For a local rewrite recognized as arithmetic, inspect the original rule
constructor. In the unary and binary cases, recognition uniqueness identifies
the requested operation and ordered operands, so the corresponding local
certificate applies. Every store, application or copying constructor has a
different source-head label. Head uniqueness therefore excludes it; binary
codes are split into addition and subtraction where needed to distinguish
assignment and atom equality. No existing stability or copyability premise is
erased. Inducting over the actual rewrite relation now derives certificates
for every arithmetic computation: the local case is just proved, while renaming,
plugging and expansion use recognition reflection and certificate transport.

\lean{Copying/ArithmeticCoherence}{107}{149}{Combined.local_arithmetic_certificate}

\lean{Copying/ArithmeticCoherence}{152}{165}{Combined.rewrite_arithmetic_certificate}

\termdef{arithmetic-successor-coherence}
Take any two actual arithmetic rewrites from the same source request. Their
derived certificates have the same head by source-head uniqueness and the same
result by the fixed recognized request. The survivor-image theorem therefore
gives a recursive graph isomorphism; their token equations give the state
isomorphism. This is \emph{arithmetic successor coherence}. It compares complete
successors produced by arbitrary contextual matches, not only the selected
redex fragments. A combined transition from rewrite mode cannot be a search
transition, so the result immediately covers the full combined step relation.

\lean{Copying/ArithmeticCoherence}{168}{181}{Combined.arithmetic_rewrite_deterministic}

\lean{Copying/ArithmeticCoherence}{183}{192}{Combined.arithmetic_step_deterministic}

An arithmetic transition consumes exactly one unit of successful cost.
For any chosen actual successor $t$, a successful source cost $n$ is equivalent
to $n=1+k$ and a successful cost $k$ from $t$. Any competing first successor
is isomorphic to $t$, and the previously proved execution-isomorphism law
transports its remaining run without changing cost. Lifting the state
isomorphism proves representation equivalence of arbitrary unbounded
successors; representation equivalence then yields observational equivalence
at every cost preorder. These results do not assume global machine determinism.

\lean{Copying/ArithmeticCoherence}{194}{197}{Combined.arithmetic_cost_spectrum}

\lean{Copying/ArithmeticCoherence}{199}{208}{Combined.arithmetic_unbounded_deterministic}

\lean{Copying/ArithmeticCoherence}{210}{214}{Combined.arithmetic_unbounded_equiv}

\Needspace{12\baselineskip}
\subsection{An exhaustive description of the remaining coherence cases}

\termdef{higher-order-head}
The historical predicate \emph{higher-order head} records either an application edge or
a term-contraction edge at the current token. At the arithmetic milestone these were
the two families awaiting complete successor comparison; application is
subsequently discharged in Section~\ref{sec:copying-application-coherence};
the predicate does not assert that either family is nondeterministic.

\lean{Copying/CombinedCoherenceFrontier}{9}{11}{Combined.HigherOrderHead}

For two actual rewrites, obtain the recognized choice of the first one
and inspect all eight choice constructors. Allocation, read, assignment and
either comparison branch use the completed store-coherence theorems. Unary and
binary arithmetic use the new theorem. The application and copy cases supply
the corresponding higher-order head. Thus every rewrite pair is coherent or
has one of these two heads. For full machine steps, a search first step has
its unique literal successor, and mixed search/rewrite steps are impossible.
The remaining case is the rewrite classification. Consequently any possibly
unresolved pair must also have rewrite mode.

\lean{Copying/CombinedCoherenceFrontier}{14}{28}{Combined.rewrite_coherent_or_higherOrder}

\lean{Copying/CombinedCoherenceFrontier}{32}{43}{Combined.step_coherent_or_higherOrder}

If the source is not a rewrite request at one of these heads, eliminate
the second alternative to obtain successor coherence outright. Applying the
one-step cost theorem gives the corresponding exact cost decomposition. The
same exhaustive disjunction lifts to the unbounded machine: fixed-depth coherent
successors become representation-equivalent successors, and the head/mode
alternative is unchanged. This identified the remaining proof obligation at the arithmetic milestone.
Section~\ref{sec:copying-application-coherence} further reduces it to copying.

\lean{Copying/CombinedCoherenceFrontier}{46}{49}{Combined.step_deterministic_away_higherOrder}

\lean{Copying/CombinedCoherenceFrontier}{52}{56}{Combined.cost_spectrum_away_higherOrder}

\lean{Copying/CombinedCoherenceFrontier}{59}{69}{Combined.unbounded_coherent_or_higherOrder}

\Needspace{12\baselineskip}
\subsection{Exact costs, arbitrary hosts and regressions}

Search in the right component of a serial graph is a genuine search step
in the whole graph, with its token inserted at the actual shared interface.
This follows from simultaneous-expansion search transport at the right-hand
component of the two-edge serial pattern.

\lean{Copying/ArithmeticCost}{8}{12}{Search.serial_right}

Unary evaluation takes three search transitions to reach its rewrite
request, one arithmetic transition, and one return from the resulting number.
Apply search cost decomposition at each search step and arithmetic cost
decomposition at the compute step. The final external-input backtrack state
has exactly zero remaining successful cost. Adding these equations gives
$n=5$; conversely the explicit remaining costs $4,3,2,1,0$ reconstruct a
successful run. Binary evaluation has five search transitions before rewriting,
then the same compute and return, giving $n=7$ in both directions. These
statements concern every successful run of the full machine.

\lean{Copying/ArithmeticCost}{19}{28}{Arithmetic.unary_cost_iff}

\lean{Copying/ArithmeticCost}{31}{42}{Arithmetic.binary_cost_iff}

For an arbitrary typed host hole, first cast the local redex boundary to
the hole's signature and plug the original arithmetic transition into that
hole. Its source head is still arithmetic by boundary transport and the match
embedding. Head uniqueness excludes both higher-order alternatives, so the
classification theorem compares any actual successor with this chosen one.
This proves unary and binary host-successor coherence while retaining the
entire host and its recursive boxes.

\lean{Copying/ArithmeticCost}{45}{58}{Arithmetic.unary_host_successor}

\lean{Copying/ArithmeticCost}{61}{74}{Arithmetic.binary_host_successor}

Exact padding reflection and exact finite/unbounded-run correspondence
transfer the local five- and seven-step characterizations to every padded
closed arithmetic redex program. The underflow regression combines the
previously constructed subtraction result $0$ for $m\leq n$ with the new
all-run cost characterization. The negation regression similarly combines
its exact result $0$ with the five-step spectrum. Neither operation is replaced
by an observationally assumed constant; the existing rule and execution proofs
supply the stated results.

\lean{Copying/ArithmeticCoherenceExamples}{9}{13}{Arithmetic.Examples.unary_pad_cost_iff}

\lean{Copying/ArithmeticCoherenceExamples}{16}{20}{Arithmetic.Examples.binary_pad_cost_iff}

\lean{Copying/ArithmeticCoherenceExamples}{23}{28}{Arithmetic.Examples.underflow_result_and_cost}

\lean{Copying/ArithmeticCoherenceExamples}{31}{36}{Arithmetic.Examples.negation_result_and_cost}

\termdef{arithmetic-cost-handoff}
The \emph{arithmetic cost handoff} retains an arbitrary host $C$ above a binary
redex. A successful cost $n$ at the redex's inserted input is equivalent to
$n=7+k$ and a successful cost $k$ from the inserted backtrack state of its
numeric result. Transport the five search steps and final numeric return
through serial composition, and use the embedded recognition and contextual
arithmetic coherence at the computation step. Summing the seven one-step
cost equations proves the forward implication; choosing the seven successive
remaining costs proves the reverse implication. The returned token is internal
to the host, so it is not mistaken for a final whole-program state. No
termination assumption is made about the host continuation.

\lean{Copying/ArithmeticCoherenceExamples}{40}{59}{Arithmetic.Examples.binary_host_handoff}

Finally consider unary negation applied to a binary arithmetic redex.
One search step enters the inner redex. Its seven-step handoff changes the
inner graph into the numeric result before the outer operation resumes.
There remain one outer readiness step, one unary computation and one numeric
return. The final external-input backtrack has cost zero. Composing these
cost equivalences proves that the nested expression succeeds at exactly
$1+7+1+1+1=11$ transitions. This regression checks actual graph change between
the two arithmetic computations, not merely concatenation of abstract numbers.

\lean{Copying/ArithmeticCoherenceExamples}{63}{81}{Arithmetic.Examples.nested_cost_iff}
